\section{High availability enabled}
\label{sec:ha:enabled}

This is the configuration a venue runs in. The switch starts a second instance
of every component that must be singular --- the sequencer, the matching engine
and the recorder --- together with the components whose only purpose is to
decide which of a pair may act: a pair of arbiters and a witness.

The switch changes nothing about the replicated peers. The authentication
services and the gateways run the same way in both configurations, for the reason
\cref{sec:ha:disabled:scope} gives: several of each run at once, none of them
leads, and an arbiter has nothing to decide about them.

\subsection{Two pairs of words, and why they must not be confused}
\label{sec:ha:enabled:vocabulary}

\textbf{Primary and secondary are permanent names.} They are fixed in
configuration when an instance is deployed and never change for the life of that
deployment. They say which instance is preferred when both start together and
neither has any other claim, and they say nothing else. They are not a statement
about which instance is doing the work.

\textbf{Leader and follower are positions.} The leader is the instance currently
entitled to act; the follower is the one that is not. Either configured instance
may hold either position, and which it holds may change repeatedly over a
deployment's life.

\textbf{``The primary is the follower'' is therefore not a contradiction.} After
one failover it is the ordinary state, and it remains so until something moves
leadership again. A design that treats the two pairs of words as synonyms moves
leadership back to the primary for no reason beyond its name, and moves it to the
instance that has just failed.

The pair must also be able to exchange positions repeatedly. If the secondary is
promoted and later dies, the primary, which is the follower by then, must be able
to take over again. A scheme that can fail over only once postpones an outage
instead of preventing it.

This document uses neither \emph{active} nor \emph{standby}. Readers take those
words to mean the configured identity in some places and the position in others,
which is the ambiguity this section exists to prevent.

\subsection{What changes, and what does not}
\label{sec:ha:enabled:what_changes}

\textbf{What changes} is that the venue survives the loss of a machine. The
configuration in \cref{sec:ha:disabled} can only restart what died and rebuild
what it held. This configuration has a second instance that already holds it.

\textbf{What does not change is the recovery itself.} For the matching engine,
becoming current after a restart and becoming current after a promotion are the
same operation. They differ only in where the starting point comes from: a file
after a restart, and a replica after a promotion. This section therefore
describes how an instance learns that it may act, and the failures that exist
only because two instances run.

\textbf{What is new} follows from a second instance existing: deciding which of
the two may act, ensuring that only one of them acts, detecting that the other
has gone, and refusing what an instance sends after its entitlement has passed to
the other. The subsections that follow cover those, and the components that exist
to arbitrate.

\subsection{What running two of everything has to be worth}
\label{sec:ha:enabled:worth}

Running two instances of a component adds ways for the venue to fail that running
one instance does not have. A single instance cannot disagree with itself, cannot
be separated from itself by a network failure, and cannot both act and fail to
act. A pair can reach all three conditions.

Two things must therefore hold for the second instance to be an improvement.

\textbf{The venue must not be stopped by something a single instance would have
survived.} If a pair of sequencers stops trading because they cannot reach each
other, when one sequencer alone would have carried on, the redundancy has made
the venue less available rather than more.

\textbf{The venue must prevent the new failures rather than make them unlikely.}
Two instances acting at once, and neither instance acting at all, are conditions
that cannot arise with one instance. They exist only because a second instance
exists. Making them rare is not sufficient, because the venue goes on looking
healthy in both. Nobody sees the first until the two instances' states are
compared, and nobody sees the second until somebody asks why nothing is trading.

Most of this section is about preventing them.

\subsection{The arbiter and the witness}
\label{sec:ha:enabled:arbiter}

Two instances of a group may act, and only one of them may act at a time. This
subsection says how the venue decides which one.

An instance cannot make this decision alone. When an instance loses contact with
its peer, it cannot distinguish a peer that has stopped from a peer it can no
longer reach. Both readings lead to a wrong action. If it promotes itself and the
peer is alive, both instances act and their state diverges. If it declines and
the peer is dead, neither instance acts and the venue stops trading. The second
failure is the safer one, and it takes longer to notice, because no process
crashes and nothing logs an error.

\textbf{A majority of three voters decides.} Each group has three voters: its two
instances and the arbiter pool. An instance acts only while a majority of the
three has granted it a lease that has not run out. Its own vote is one of the
three, so it needs a grant from one other voter: its peer, or the arbiter pool. A
voter grants a lease to one instance at a time, and promises not to grant another
until the lease period has passed. Two instances would each need two of the three
votes at the same moment, so at least one voter would have to have granted both,
and no voter does. Every message also carries the generation number of the
entitlement it was sent under, and a receiver refuses a message whose generation
has been superseded.

\textbf{Two arbiters run, and one witness runs alongside them.} The arbiter pool
votes through whichever arbiter is active. The two arbiters decide which of them
is active by the same rule, with the witness as their third voter: an arbiter is
active only while a majority of the two arbiters and the witness has granted it a
lease. Losing one arbiter moves no entitlement and interrupts nothing, which is
\reqref{R-0063}. The witness has no pair and no standby.
\Cref{sec:ha:enabled:arbiter:witness} states what the venue does about that.

\subsubsection{What is done about the witness being single}
\label{sec:ha:enabled:arbiter:witness}

\textbf{The witness holds no state, so losing it loses nothing.} It has no record
to recover and no history to lose. After it restarts it grants nothing for one
lease period, which is the whole of its recovery.

\textbf{Losing the witness alone changes nothing.} Two arbiters that can see each
other grant each other leases without it. The absence of the witness is felt only
when an arbiter is also lost: the surviving arbiter then holds only its own vote,
and it is active only while its existing lease lasts. The pool then has no active
arbiter until another voter returns. Component leaders are unaffected while their
peers are running, because each renews its lease with its peer.

\textbf{The witness is off every ordinary path.} Only the arbiters talk to it. No
order, report or component touches the witness.

\textbf{The witness runs where the arbiters do not.} A witness that shares a
failure domain with one arbiter gives little protection, because a single event
then removes that arbiter and the witness together, and the surviving arbiter
holds only its own vote. The witness therefore runs on a third machine, on a
different power supply, behind a different network switch, and ideally in a
different rack. Its usefulness comes from where it is deployed rather than from
what it does.

Three properties of the witness process make that placement cheap:

\begin{itemize}[nosep]
  \item \textbf{It needs no storage.} It holds nothing on disk, so it depends on
        no array and no filesystem.
  \item \textbf{It is off the hot path.} It needs no pinned cores, no low-latency
        hardware and no proximity to any other component, so it can run on the
        least capable machine in the deployment.
  \item \textbf{Its network route matters more than its speed.} The failure that
        hurts the venue is an arbiter being unable to reach the witness. The
        route to the witness should therefore fail independently of the route
        between the two arbiters.
\end{itemize}

The venue goes on trading when both arbiters and the witness are gone, for as long
as each group's leader can renew its lease with its peer. What it cannot do then is
replace a leader that fails: \cref{sec:ha:enabled:degraded} says what happens.

\subsubsection{What state it holds}
\label{sec:ha:enabled:arbiter:state}

The arbiter holds whether it is active and, while it is, one vote per group: the
instance it has promised its vote to, if any, and until when. It also holds the
highest generation granted in each group. It holds nothing about orders, sessions
or sequence numbers. That limit is deliberate: a component that decides which
instance may act must not itself hold state that can be wrong.

The witness holds its vote on which arbiter is active, and nothing else.

\subsubsection{Where the copy that outlives it lives}
\label{sec:ha:enabled:arbiter:durable}

The arbiter keeps nothing across its own restart. It does not need to know what it
promised before it stopped: it grants nothing for one lease period after starting,
by which time every such promise has run out. The same applies to an arbiter that
becomes active in place of the other, because it does not know what the other
promised.

The highest generation granted in each group is the one thing worth keeping, and
the active arbiter tells the passive one each time it rises. An arbiter that
becomes active therefore knows it, unless both arbiters have restarted.

\subsubsection{How a restarted instance becomes current}
\label{sec:ha:enabled:arbiter:recovery}

A restarted arbiter waits one lease period and then votes. Until then it refuses
every request, which is \reqref{R-0059}.

\subsubsection{What the member is told}
\label{sec:ha:enabled:arbiter:member}

Nothing. The arbiter sits off the path a member's orders travel, and a member
cannot tell that one exists. An arbiter's failures reach a member only through
the components that consult it, so the requirements here are stated in terms of
those components rather than of members.

\begin{requirement}{R-0058}{Two instances of a group are never both entitled to act}
  At any moment, at most one instance of a group shall hold the entitlement to
  act.
  \rationale{This is the property the arbiter exists to provide, and the rest of
  this subsection follows from it. Two instances acting on the same state
  diverge, and nothing afterwards can establish which of them was right.}
  \uncovered{no scenario asserts the property directly, only particular routes to it}
\end{requirement}

\begin{requirement}{R-0146}{An instance acts only while a majority of its three voters grants it a lease}
  An instance of a group shall act only while it holds a lease, not yet run out,
  from its peer or from the arbiter pool. It shall count each lease from the
  moment it asked for it, less an allowance for clocks that run at slightly
  different rates. A voter shall grant a lease to one instance at a time, shall
  grant nothing for one lease period after it starts unless it has kept a record,
  made during the same boot of its machine, of what it promised, and shall grant
  no generation below one it has already granted. An instance that has granted its
  peer a lease shall not ask to lead until that lease has run out.
  \rationale{An instance that loses contact with its peer cannot tell a dead peer
  from an unreachable one, so no rule it applies alone can be right in both
  cases. A majority of three can be: two instances would each need two of the
  three votes at the same moment, which would need a voter that has granted both.

  Each clause closes a way in which two instances come to act at once, and model
  checking showed each one is needed: without it, two instances act as leader at
  once within a few steps. Counting the lease from the request rather than the
  grant means the holder always believes its lease ends first. Waiting after a
  start covers a voter that has forgotten what it promised; a voter that kept a
  record has not forgotten, and so can resume at once, which is what lets a
  supervised restart keep the lead. The epoch rule stops
  a voter granting an older generation. An instance that granted its peer a lease
  must not then count its own vote for itself.}
  \covers{ha_test:8, ha_test:15, ha_test:39, ha_test:49, tla:lease-1-holder-counts-from-arrival, tla:lease-2-candidate-ignores-own-grant,
  tla:lease-3-restart-does-not-wait}
\end{requirement}

\begin{requirement}{R-0059}{An arbiter that does not know what was promised grants nothing until every promise has run out}
  An arbiter that has just started, or has just become the active arbiter, shall
  grant no lease to any component for one lease period.
  \rationale{An arbiter keeps nothing across its own restart, and an arbiter that
  becomes active does not know what the previously active one promised. During
  that period it cannot distinguish two situations: a group in which no instance
  holds a lease, and a group whose leader holds one it cannot see. Granting a
  lease is correct in the first situation. In the second it gives a group two
  instances entitled to act, which is the outcome the arbiter exists to prevent.

  Waiting one lease period settles it without any knowledge: by then every
  promise made before is certain to have run out. The leader of each group
  renews with its peer meanwhile, so the wait costs nothing while a pair has both
  instances running.}
  \uncovered{no scenario asks a freshly started arbiter about a group it has not been told of}
\end{requirement}

\begin{requirement}{R-0060}{Silence from an arbiter is not read as the absence of one}
  A component that asks and hears nothing shall ask again before concluding that
  no arbiter exists.
  \rationale{An arbiter that stays silent may be the passive one, or one waiting
  out a lease period under \reqref{R-0059}. A component that treats the silence as
  an absent arbiter and acts anyway creates the second leader the silence was
  preventing. A component therefore leads only on a grant, never on the absence
  of an answer.}
  \uncovered{no scenario has a component hear nothing from any voter and checks that it does not lead}
\end{requirement}

\begin{requirement}{R-0061}{Entitlement is held under a lease and lapses unless renewed}
  An instance's entitlement to act shall expire unless it is renewed, and an
  instance whose entitlement has expired shall stop acting at once.
  \rationale{A record that never expires states only that something was true at
  the moment it was written. Presence alone cannot be the test either: an
  instance that holds its connection open and stops renewing looks, from the
  outside, the same as a healthy one. An instance that keeps acting after its
  lease has run out is a second leader the moment a voter grants its peer a
  lease, so stopping is what makes the lease mean anything.}
  \uncovered{no scenario stops a leader renewing while it stays connected}
\end{requirement}

\begin{requirement}{R-0099}{Adjusting a machine's clock does not move an entitlement}
  The expiry of a lease shall be measured by a clock that only moves forward at a
  steady rate, and shall not be affected by the machine's notion of the time of
  day being corrected, stepped or set.
  \rationale{Time-of-day clocks are adjusted routinely, by a time service
  correcting drift, by a leap second, or by an operator. A lease measured against
  such a clock can expire while it is being renewed properly, which moves an
  entitlement for a reason unconnected to any component's health. An adjustment
  in the other direction is worse: a lease that should have expired goes on being
  honoured after its holder has gone.}
  \uncovered{no scenario adjusts a clock}
\end{requirement}

\begin{requirement}{R-0109}{An instance that is not making progress does not keep its entitlement}
  Holding an entitlement shall depend on the instance continuing to do the work
  the entitlement is for, and not only on its continuing to ask to keep it.
  \rationale{Renewing an entitlement requires no more than a running thread. An
  instance starved of processor time, or stalled in a way that leaves its timers
  running, renews on schedule and does nothing else. Neither presence nor renewal
  reports whether the instance is doing any work, so both report such an instance
  as healthy. The venue then stays with an instance that has stopped, and
  everything waits while nothing appears to be wrong.}
  \uncovered{no scenario starves an instance that keeps renewing}
\end{requirement}

\begin{requirement}{R-0115}{How long an entitlement lasts without renewal is a stated figure}
  The period after which an unrenewed entitlement lapses shall be a figure the
  venue states, and it shall be long enough that an instance doing its work is
  never mistaken for one that has stopped.
  \rationale{A period that is too short takes the entitlement away from a busy
  instance for being busy, which moves the venue's work to the other instance at
  the moment there is most of it. A period that is too long leaves the
  entitlement with an instance that has genuinely stopped, and nothing acts until
  it lapses. The figure alone shows neither error, so it has to be derived from
  how long the work takes.}
  \uncovered{no scenario tests the period, and none is stated to test}
\end{requirement}

\begin{requirement}{R-0062}{An instance already acting keeps its entitlement when another returns}
  Where an instance holds the entitlement and is still present, an instance
  rejoining the group shall be told to follow, whatever its configured name.
  \rationale{Preferring the configured primary at this point moves the
  entitlement for no reason but a name. It also moves the entitlement to the
  instance that has just failed, and away from one that is serving. The
  configured preference is a cold-start tie-break.}
  \covers{ha_test:24, ha_test:25}
\end{requirement}

\begin{requirement}{R-0063}{Losing one arbiter, or the witness, changes nothing}
  The failure of a single arbiter, or of the witness, shall not move any
  group's entitlement and shall not interrupt the venue.
  \rationale{Redundancy is what makes the arbiter available, so nothing needs to
  watch it. A design in which losing one arbiter mattered would need something to
  supervise the arbiter, and that supervisor would need supervising in turn.}
  \covers{ha_test:2, ha_test:4, ha_test:5, ha_test:33}
\end{requirement}

\begin{requirement}{R-0064}{A generation is never issued twice, and never goes backwards}
  Each grant of entitlement shall carry a generation greater than every
  generation the venue has used. Where every instance that could have issued one
  has restarted, the new generation shall still exceed every generation
  previously used, which requires that the number outlive the processes rather
  than any one of them.
  \rationale{The generation is what lets a receiver refuse a superseded sender. A
  generation held only in memory restarts at the beginning, so a pair that
  restarts together agrees a generation the venue has used before, and a message
  left over from the earlier one then compares as current.

  This is the requirement most exposed to a restart of everything at once. Where
  each instance keeps its own record, a group whose members all restart can still
  exceed every generation used, because each instance reads its own record and
  the highest wins. Where the record is kept in one place, the survival of that
  place is the whole of the guarantee. This chapter does not settle which of the
  two the venue uses. The requirement holds either way, and one of them has to be
  chosen deliberately.}
  \uncovered{no scenario restarts every instance of a group together}
\end{requirement}

\begin{gap}{}
  A generation can go backwards. An arbiter that has restarted, and has not been
  told the highest generation by the other arbiter because it has restarted too,
  can grant a new leader a generation below one already led in. Two instances
  still never act at once, but receivers that saw the higher generation refuse
  what the new leader sends. The new leader learns of the higher generation when
  its peer refuses it, and asks again above it. Where the peer is also down, only
  a receiver could tell it, and receivers do not yet do so. See
  \texttt{docs/availability/tla/findings.md}, section 11.5.
\end{gap}

\begin{requirement}{R-0065}{An instance whose entitlement has passed has what it sends refused}
  A receiver shall refuse a message carrying a generation older than the newest
  it has accepted for that group.
  \rationale{An instance whose lease has run out stops acting by itself
  (\reqref{R-0146}), before anything it sends could be refused. Refusing an older
  generation is the second defence, for anything that reaches a receiver all the
  same: a message delayed in transit, or a leader granted a generation below one a
  receiver has already seen. It needs no special hardware and no cooperation from
  the instance being refused.}
  \uncovered{a superseded instance stops before it sends, so no scenario makes one send anything for a receiver to refuse}
\end{requirement}

\subsection{The sequencer pair}
\label{sec:ha:enabled:sequencer}

Two sequencers run. One of them holds the entitlement to act. The other receives
everything the acting instance does, and waits.

\subsubsection{What state it holds}
\label{sec:ha:enabled:sequencer:state}

The pair holds the same state as the single sequencer in
\cref{sec:ha:disabled:sequencer}: the ordering, the log, the generation and each
session's numbering. The pair holds it twice, on two machines with two disks.

The follower builds its copy from the same records, in the same order, as the
leader writes them. A session's numbering reaches the follower by the same route
it reaches the leader, because gateways report to both instances. The follower
therefore holds what a restarted instance would have to be told.

\subsubsection{Where the copy that outlives it lives}
\label{sec:ha:enabled:sequencer:durable}

The copy lives on the other machine, which is the point of this configuration. A
failed disk, a stopped machine and a hung kernel each leave the venue's record of
what it has done intact, because the venue wrote a second copy as it wrote the
first.

\textbf{The rule that makes the second machine worth running governs when a
member is told.} An order becomes durable in two stages. The leader writes the
order, which permits the venue to give the order to the matching engine. The
follower then acknowledges the order, which permits the venue to give the
resulting execution report to the gateway. A member is therefore never told about
an order that exists on one machine only.

That ordering adds a round trip to the reporting path, and it is what gives the
venue the property the second machine is run for. Removing the round trip to
shorten the path would remove that property.

\subsubsection{How a restarted instance becomes current}
\label{sec:ha:enabled:sequencer:recovery}

A promoted follower is already current, or close to it. It holds the records and
has acknowledged them, and it needs only to apply what it has received and not
yet applied.

\textbf{A promotion recovers more than the restart in
\cref{sec:ha:disabled:sequencer} does.} A restarted sequencer recovers its
sequence number from its log, its generation from a file, and its sessions by
being told, which leaves it holding nothing for a session that was not bound at
the time. A promoted follower was told the same things as they happened, so it
holds the numbering of sessions that no gateway would report.

The pair also settles leadership between themselves when they can see each other,
rather than asking the arbiter. Both instances hold both instance identities and
both generations, so both compute the same answer from the same inputs. The
arbiter decides the case the pair cannot settle alone: one instance cannot see
the other, and a second claimant may exist that it cannot see.

\subsubsection{What the member is told}
\label{sec:ha:enabled:sequencer:member}

Nothing. The member sees a brief pause in responses and no break in numbering.

\begin{requirement}{R-0104}{A member is told nothing false while an entitlement is moving}
  Between the loss of the instance that was acting and the moment its replacement
  has caught up, a member may receive no answer to what it sends, and shall not
  receive an answer that is wrong. An order shall not be rejected as unknown, and
  an order shall not be reported accepted, on the basis of state that is not yet
  current.
  \rationale{The interval is short, but it is not zero, and a member cannot tell that it is in
  one. Silence during the interval is a delay, and a member is built to tolerate a
  delay. An answer given from state that is still being assembled is a statement
  the venue cannot support, and the member acts on that statement.}
  \uncovered{no scenario sends orders during a failover and checks the answers}
\end{requirement}

\begin{requirement}{R-0066}{A member is never told of an order held on one machine only}
  An execution report shall not be sent to a member until the record of the order
  it concerns has been acknowledged by a second instance.
  \rationale{An order that exists in one place is lost when that place is lost, so telling a
  member about such an order commits the venue to something it may not be able to
  honour. This is the property the second machine is run for, and the other
  requirements in this section depend on it.}
  \uncovered{no scenario asserts that a report waits for the follower}
\end{requirement}

\begin{requirement}{R-0067}{Losing a machine does not lose a committed order}
  An order committed before a machine was lost shall be processed and answered by
  the surviving instance.
  \rationale{Committing the order is the venue taking it. Losing one of the two machines
  holding that order is the failure the pair exists to absorb.}
  \uncovered{no scenario asserts an order committed before the kill is answered after it}
\end{requirement}

\begin{requirement}{R-0068}{The venue keeps trading when the acting sequencer is lost}
  Members shall continue to be able to place orders across the loss of the
  sequencer instance that was acting, without reconnecting and without being
  told to retry.
  \rationale{This is the most visible reason for running two machines, and a member would
  notice immediately if it did not hold. It is stated separately from the
  requirements about what survives because continuing to trade and keeping what was
  taken are different properties. A venue can meet this requirement while failing
  every other requirement in this section.}
  \covers{ha_test:1}
\end{requirement}

\begin{requirement}{R-0069}{A member's session survives a failover}
  A member's numbering, and the orders open under its session, shall be
  unaffected by leadership moving between sequencers.
  \rationale{A failover the member can detect has not achieved what it was for. The venue
  tells the follower each session's position as that position advances, for this
  reason.}
  \uncovered{no scenario fails the sequencer over under a live session}
\end{requirement}

\begin{requirement}{R-0070}{A follower that falls behind stops the venue confirming, rather than being dropped}
  Where the second instance cannot keep up, the venue shall stop sending
  execution reports rather than continue without it.
  \rationale{Dropping the follower to keep the reporting path moving abandons \reqref{R-0066}
  without reporting it, at the moment that requirement matters most. Stopping the
  reports instead leaves members seeing orders received and no reports, which a
  member can see and recover from. A member cannot see the venue running on one
  machine.}
  \uncovered{no scenario makes the follower lag}
\end{requirement}

\begin{requirement}{R-0100}{Waiting for a disk or for a peer does not stop a component working}
  A component shall continue to process what arrives while it is waiting for a
  write to complete or for a second instance to acknowledge one, and shall
  continue to renew its entitlement and to answer enquiries about its health
  throughout.
  \rationale{The requirements in this section make a component wait for a disk and for a peer,
  and both can be slow for reasons unconnected to the venue. A component that stops
  processing while it waits also stops answering members, stops replying to its
  peer and stops renewing its lease. From outside, such a component looks the same
  as one that has stopped, so the venue restarts it, and a slow disk becomes a
  restart rather than a delay.}
  \uncovered{no scenario makes a disk or a peer slow}
\end{requirement}

\begin{requirement}{R-0110}{How far behind a second instance may fall is a stated figure}
  The point at which a second instance is too far behind shall be a figure the
  venue states, expressed in what has not been acknowledged rather than in
  elapsed time alone.
  \rationale{\reqref{R-0070} requires the venue to stop confirming rather than proceed without
  the second instance. Neither implementing nor testing that requirement is possible
  without knowing when the point is reached. A figure expressed only as a duration
  moves with the rate orders arrive, so a venue under load reaches it after far
  fewer records than a quiet venue does.}
  \uncovered{no scenario makes a second instance lag, and no figure says when it is too far}
\end{requirement}

\begin{requirement}{R-0071}{A sequencer that finds its own record damaged stops rather than serving from it}
  An instance that cannot read its log whole shall not act on the part it can
  read, and the instance holding an undamaged copy shall act instead.
  \rationale{This is the answer that \cref{sec:ha:disabled:recovery} states the disabled
  configuration does not have. An instance serving from a partial log assigns
  numbers it has assigned before, and answers orders other than the ones it
  committed.}
  \uncovered{no scenario damages a log}
\end{requirement}

\begin{requirement}{R-0072}{A pair that can see each other settles leadership between themselves}
  Where both instances are present and neither holds the entitlement, they shall
  reach the same answer without referring it, and shall reach it whichever
  started first.
  \rationale{Both instances hold both identities and both generations, so both compute the
  same result from the same inputs. Asking the arbiter instead makes the venue wait
  on a component that has just started itself, and that holds less information than
  either instance.}
  \covers{ha_test:36, ha_test:37}
\end{requirement}

\subsection{The matching engine pair}
\label{sec:ha:enabled:me}

Two matching engines run. One of them acts. The other holds a replica of the set
of open orders, built from the same records in the same order.

\subsubsection{What state it holds}
\label{sec:ha:enabled:me:state}

The pair holds the set of open orders and the sequence number that set is current
to, as \cref{sec:ha:disabled:me:state} describes, and holds them twice. The
engine keeps no log of its own in either configuration. The sequencer holds the
record of what was taken, and what the engine holds follows from that record.

\subsubsection{Where the copy that outlives it lives}
\label{sec:ha:enabled:me:durable}

The copy lives in the other instance, kept current by the same stream.
\Cref{sec:ha:disabled:me:durable} states that this copy does not exist when high
availability is disabled, which is why that configuration needs a checkpoint.

\subsubsection{How a restarted instance becomes current}
\label{sec:ha:enabled:me:recovery}

The promoted instance reports the sequence number its replica has reached,
receives everything committed after that number, applies it silently, and begins
acting.

\textbf{This is the same operation as \cref{sec:ha:disabled:me:recovery}
describes.} A restarting instance learns its starting number from a checkpoint,
and a promoted instance takes it from the replica it has been maintaining.
Nothing else about the operation differs, so the mechanism built for the disabled
configuration is the mechanism this configuration uses.

The engine differs from the sequencer in one respect. A sequencer pair settles
leadership between themselves in both directions, because both instances hold
both generations at the same moment. A matching engine only ever \emph{gives}
leadership up. It adopts the follower position on its peer's assertion, takes the
generation it is given, and issues no generation of its own. An instance that
gives up leadership cannot produce a second acting instance, whatever its peer
has got wrong, which is why acting on the peer's assertion alone is safe here.

\subsubsection{What the member is told}
\label{sec:ha:enabled:me:member}

\begin{requirement}{R-0073}{An open order survives a failover}
  An order open before leadership moved shall be open after it, on the same
  terms, and shall remain cancellable by the member that placed it. Where the set
  cannot be recovered, \reqref{R-0020} applies instead and every order the venue
  can still name is cancelled to its member.
  \rationale{\reqref{R-0018} requires this of a restart. The configuration with high
  availability enabled exists to give a member a better outcome, so a member must
  not keep its orders when a process dies and lose them when a machine dies.

  The two requirements form an order of preference rather than a choice. The venue
  recovers what it can recover, and reports what it cannot. The member is told
  exactly where it stands in either case, and the venue cancels only where it does
  not have the orders.}
  \uncovered{scenario 21 asserts the fallback, and nothing asserts this}
\end{requirement}

\begin{gap}{}
  \textbf{The venue does not meet this requirement as it is configured, and says
  so.} What a promoted matching engine does with the orders it inherits is a
  stated policy, \texttt{order\_book.open\_orders\_on\_promotion}, and it is set
  to \texttt{cancel} in every environment. Every promotion therefore takes the
  fallback path: each order is cancelled, its member is told, and the venue leads
  an empty book.

  The engine can give the other answer, because \texttt{keep} carries the book
  across. The mechanism exists, and what is missing is the grounds for trusting it.
  A promoted instance assembles its book from two sources, and \textbf{one of them
  is checked and the other is not}. The catch-up applied before the instance acts
  is accounted for, which is \reqref{R-0101}. The replica the instance maintained
  by \texttt{BookUpdate} while it was the follower is not accounted for: it applies
  those updates as they arrive, and nothing establishes that none was missed.
  Keeping a book that is short of an order is worse than cancelling every order,
  because the member is told nothing and believes it holds an order that the venue
  does not.

  The setting therefore stands at the outcome that is worse and certain. Changing
  it requires the same check to be applied to the replication stream that
  \reqref{R-0101} applies to the catch-up.
\end{gap}

\begin{requirement}{R-0074}{An instance that gives up leadership cannot thereby create a second one}
  An instance shall adopt the follower position only on its peer's assertion that
  the peer leads, only while itself holding no position, and only where the
  generation offered is not older than one it has already seen.
  \rationale{An instance that gives up a position cannot produce two acting instances,
  whatever its peer has got wrong. A one-way announcement is therefore sufficient
  here, where the sequencer needs an exchange.}
  \uncovered{no scenario offers a stale generation to an instance holding no position}
\end{requirement}

\begin{requirement}{R-0075}{The venue keeps working when the acting engine is lost}
  Members shall continue to be able to place and cancel orders across the loss
  of the matching engine instance that was acting.
  \rationale{Cancelling matters here as much as placing. A member that can place orders but
  cannot cancel the orders it already holds is worse off than a member that can do
  neither, because it accumulates exposure that it has no way to reduce.}
  \covers{ha_test:16}
\end{requirement}

\begin{requirement}{R-0076}{Orders taken while no engine was acting are applied on promotion}
  An order committed while no instance was acting shall be applied and answered
  once one is.
  \rationale{The venue took those orders, so it must answer each of them. The disabled
  configuration cannot serve this case, because it has no second instance holding a
  position to catch up from.}
  \uncovered{no scenario asserts that deferred orders are answered after a promotion}
\end{requirement}

\subsection{The publisher and recorder pairs}
\label{sec:ha:enabled:publishing}

Two publishers run, and two recorders.

\subsubsection{What state it holds}
\label{sec:ha:enabled:publishing:state}

The publisher holds what it has been given, and holds no position on behalf of
any subscriber. A subscriber presents its own position when it connects, as
\cref{sec:ha:disabled:mep:state} describes. The recorder holds one number: how
far it has published and had confirmed.

\subsubsection{Where the copy that outlives it lives}
\label{sec:ha:enabled:publishing:durable}

Both instances of a publisher receive the same records, identified by the same
sequence numbers. Those numbers are the sequencer's, and the publisher assigns
none of its own. \textbf{A subscriber's position therefore means the same thing
at either instance}, so a subscriber moving to the surviving instance needs no
translation, and the two publishers need no handover between them.

The recorder's position needs less still. The recorder writes it only after the
bus confirms the publish, so the position is a lower bound on what has been
published. A promoted instance resuming from a stale copy publishes some records
again and loses none, and the consumers downstream are required to handle the
duplicates.

\subsubsection{How a restarted instance becomes current}
\label{sec:ha:enabled:publishing:recovery}

Other components tell both of them. Each subscriber tells a publisher where that
subscriber had reached. A recorder reads its own position and presents it to a
publisher. Neither needs a handover from the instance it replaced, so neither
pair needs the coupling that the sequencer pair has.

\subsubsection{What the member is told}
\label{sec:ha:enabled:publishing:member}

The publisher tells a member nothing. A failure of the recorder reaches members
only through \reqref{R-0053}, and running two recorders is what keeps it from
reaching them at all.

\begin{requirement}{R-0077}{A subscriber resumes at the surviving instance without loss}
  Where the instance a subscriber was reading from is lost, the venue shall
  supply every record after the position that subscriber presents, from the
  instance it moves to.
  \rationale{A subscriber's position is expressed in the venue's own sequence numbers rather
  than in numbers a publisher assigns, so either instance can act on it. The two
  publishers hand nothing over.}
  \uncovered{no scenario moves a subscriber between publisher instances}
\end{requirement}

\begin{requirement}{R-0078}{A promoted recorder republishes rather than skips}
  Where publishing to the bus is taken over by another instance, the venue shall
  resume from a position no later than the last publish the bus confirmed, and
  shall publish again whatever it cannot establish was confirmed.
  \rationale{The position is a lower bound, so resuming from a stale position costs duplicate
  records, and resuming from a position later than the confirmed one costs records
  entirely. \reqref{R-0049} permits the first and forbids the second.}
  \uncovered{no scenario, and none is possible until the recorder exists}
\end{requirement}

\begin{requirement}{R-0079}{Losing one publisher or one recorder does not halt the venue}
  The loss of a single publisher instance, or of a single recorder instance,
  shall not stop members trading.
  \rationale{\reqref{R-0053} halts the venue when no recorder is publishing. Running two
  recorders is what makes that condition rare. Halting on the loss of one of them
  would produce the outage that the second recorder exists to prevent.}
  \uncovered{no scenario, and none is possible until the recorder exists}
\end{requirement}

\subsection{A process dies while its peer is present}
\label{sec:ha:enabled:process_death}

A component dies on a machine that is otherwise healthy, and the launcher
restarts it, while a second instance of the same component runs elsewhere. Both
configurations apply to this case.

\subsubsection{What is at risk}
\label{sec:ha:enabled:process_death:risk}

The venue's ability to trade is not at risk, because the peer can act and trading
continues whatever the venue decides. What is at risk is \textbf{the pair}.

If the peer promotes, one instance serves the venue until somebody restores the
other, and the venue is exposed to a second failure for as long as that takes. If
the launcher restarts the dead instance instead, the pair is whole again in
seconds and the venue loses nothing.

Nothing the venue does can restore a machine that has died, so promoting the peer
is the only remedy for that failure. The venue can restore a process that has
died on a living machine, so promoting the peer in that case gives up a recovery
that was available. The venue must therefore treat the two failures differently.

\subsubsection{What the venue must do}
\label{sec:ha:enabled:process_death:behaviour}

\begin{requirement}{R-0080}{A process death on a healthy machine does not cost the pair}
  Where a component dies and is restarted on a machine that has not failed, the
  venue shall end with two instances running, as it began.
  \rationale{The measure of this requirement is whether the pair is whole again, rather than
  whether the venue kept trading. A promotion keeps the venue trading on its own,
  and leaves one instance down with nothing watching for a second failure.}
  \uncovered{no scenario asserts the pair is whole afterwards, only that trading continued}
\end{requirement}

\begin{requirement}{R-0081}{The peer does not promote while a restart could still succeed}
  An instance that has lost contact with the acting instance shall wait, before
  seeking to act, for at least as long as a supervised restart takes.
  \rationale{A shorter wait fails the venue over for a failure that did not require it. A much
  longer wait tolerates a genuinely dead machine for longer than it should. The
  period is therefore the time a restart takes plus a margin. That time is not
  known until restarts have been timed, so the period is measured rather than
  chosen, and it must be measured again whenever the components change.
  The peer cannot take over until its own promise to the acting instance and the
  arbiter's have both run out, which is at least one lease period after the last
  renewal. A restarted instance that kept its promises on disk asks again as soon
  as it starts, so a restart that completes within the lease period keeps the
  lead.}
  \covers{ha_test:26, ha_test:27, ha_test:28}
\end{requirement}

\begin{requirement}{R-0082}{An instance that keeps failing does not disturb a healthy peer}
  A restarted instance shall not take the entitlement to act from a peer that
  holds it and is present.
  \rationale{Without this requirement, an instance that fails repeatedly takes leadership on
  each restart, dies, and hands leadership back, interrupting a peer that was
  working correctly once per failure. \reqref{R-0062} states the arbiter's side of
  this rule, and this requirement states the instance's side.}
  \covers{ha_test:29, ha_test:30, ha_test:31}
\end{requirement}

\begin{requirement}{R-0125}{An instance that interrupts the venue twice is not restarted again}
  Where restarting an instance has already cost the venue one interruption, a
  second interruption from that instance shall leave it stopped rather than
  restarted.
  \rationale{An interruption is a period in which a member can neither place an order nor
  withdraw one, and is told nothing while it lasts. The venue can account for one
  interruption: a component failed, and recovery took as long as it took. A second
  interruption from the same instance shows that the instance is unfit to serve,
  and restarting it again produces a third.

  Stopping the instance is what allows the pair to do its work. An instance that
  the supervisor keeps restarting returns before the arbiter can grant the
  entitlement to its peer, so the peer never acts. The venue then takes an
  unbounded series of interruptions instead of the single failover the pair was
  built for. Leaving the failing instance stopped produces the outcome the member
  should have had at the first fault.

  \textbf{The count is taken over a session rather than over a period.} A time
  window cannot be explained to the member it affects, because it makes two
  failures count together when they fall close in time and separately when they do
  not. The requirement is stated of interruptions rather than of failed starts for a
  related reason: an instance that dies every few minutes starts correctly every
  time.}
  \uncovered{no scenario flaps an instance and asserts it is left stopped}
\end{requirement}

\begin{requirement}{R-0126}{The venue does not run with no instance of a component without saying so}
  Where the last instance of a component is stopped and not restarted, the venue
  shall report that it cannot perform that component's work, rather than continuing
  in a state where nothing does it.
  \rationale{Both instances may share the cause of the failure: an input that kills whatever
  reads it, a build, or a configuration both instances load. The peer then meets
  that cause in its turn and is stopped in its turn. The venue has no instance of
  that component left, and nothing is trying to produce one.

  This failure looks like health. Every process that remains is running normally,
  no alert names a component, and only somebody who counts the processes can see
  what is missing. Reporting the condition costs an announcement. Not reporting it
  costs the time between the last instance stopping and somebody noticing, and that
  time is unbounded.}
  \uncovered{no scenario stops every instance of a component and asserts the venue says so}
\end{requirement}

\subsubsection{When the restart and the promotion race}
\label{sec:ha:enabled:process_death:race}

The two remedies run at the same time, and neither detects the other. The
launcher restarts the instance that died while the peer counts down to asking
whether it may act. \textbf{Exactly one instance ends up acting, whichever of the
two happens first.}

\textbf{The restart completes first.} The returning instance holds no position
and asks which position it holds. No instance holds the entitlement, so the
arbiter grants it to the returning instance, which resumes acting. The peer is
still waiting, finds the entitlement held, and remains the follower. The
entitlement did not move, and the pair is whole.

\textbf{The peer asks first.} The arbiter grants the peer the entitlement, and
the peer begins acting. The returning instance asks, finds the entitlement held
by an instance that is present, and becomes the follower. Neither its configured
name nor its having held the entitlement minutes earlier affects that answer.

\textbf{Both ask at once.} One party answers both claims. The arbiter grants the
entitlement to one instance, records what it granted, and refuses the other claim
against that record.

In every ordering, exactly one instance ends up acting, and no instance has to
detect what the other was doing. That is what makes the outcome certain rather
than merely likely.

\begin{requirement}{R-0090}{A returning instance assumes nothing and asks}
  An instance that has restarted shall hold no position until it is told which
  it holds, whatever position it held before and whatever it is configured as.
  \rationale{An instance that assumes it leads, because it led before or because of its
  configured name, becomes a second acting instance whenever the peer has already
  been promoted. Nothing has to go wrong with the peer for that to happen.}
  \covers{ha_test:24, ha_test:25}
\end{requirement}

\begin{requirement}{R-0091}{A returning instance that finds the position vacant takes it}
  Where an instance returns and no instance holds the entitlement, it shall take
  it rather than wait.
  \rationale{A returning instance holds no position and asks, which is what stops a restart
  producing a second acting instance. This requirement states the other half of
  that rule. Where no instance holds the entitlement, the returning instance must
  take it. Without this, a restart that completed before its peer's promotion would
  leave both instances following and none acting.}
  \covers{ha_test:26, ha_test:27, ha_test:28}
\end{requirement}

\begin{requirement}{R-0092}{Simultaneous claims are resolved by one party, not by the claimants}
  Where a returning instance and its peer both seek the entitlement at once, at
  most one shall be granted it, and neither shall be required to have detected
  the other.
  \rationale{Two instances that must detect each other in order to stay out of each other's
  way will eventually fail to do so, and they will fail at the moment they are
  least able to communicate. Referring both claims to one party removes the need to
  detect anything.}
  \uncovered{no scenario times a restart to collide with a promotion}
\end{requirement}

\subsection{Losing a machine}
\label{sec:ha:enabled:machine_loss}

A machine stops. This is the failure the configuration exists for, and it is
always a combined failure, because a machine hosts several components and losing
it loses all of them at the same instant.

\subsubsection{What is at risk}
\label{sec:ha:enabled:machine_loss:risk}

Every component that machine was running is at risk at once, together with every
position any of them held. Nothing on that machine can be restarted, so for each
of those groups the only remedy is the instance on the other machine.

Two requirements follow that do not arise when a single process dies.
\textbf{The surviving machine must be able to carry the whole venue}, rather than
the share it was carrying. \textbf{The components that decide which instance may
act must not all have been on the machine that stopped}, or the venue loses the
ability to decide at the moment it most needs to.

\subsubsection{What the venue must do}
\label{sec:ha:enabled:machine_loss:behaviour}

\begin{requirement}{R-0083}{The venue survives the loss of any one machine}
  Members shall continue to trade across the loss of any single machine,
  whichever machine it is and whatever it was running.
  \rationale{This requirement is stated of any machine rather than of a nominated one. A
  deployment that survives losing the secondary and not the primary does not meet
  it, and nothing distinguishes the two deployments until the wrong machine fails.}
  \covers{ha_test:48}
\end{requirement}

\begin{requirement}{R-0084}{No single machine holds enough of the venue to prevent a decision}
  The venue shall be deployed so that the loss of any one machine leaves it
  still able to decide which instance may act.
  \rationale{An arbiter pair and its witness deployed on one machine share that machine's
  failure. Losing the machine then leaves the venue with no arbitration at all.}
  \uncovered{no scenario deploys instances badly, and nothing would report it}
\end{requirement}

\begin{requirement}{R-0111}{A deployment that could not survive losing a machine is reported when it starts}
  Where instances are placed such that the loss of one machine would leave the
  venue unable to decide which instance may act, that shall be reported when the
  venue starts rather than discovered when the machine is lost.
  \rationale{\reqref{R-0084} is a property of where the processes were placed. A deployment
  that breaks it behaves the same as one that does not, until the wrong machine
  fails. The venue can check the placement when the instances announce themselves,
  and that is the only moment before the failure at which anybody would look.}
  \uncovered{no scenario deploys instances badly, and nothing would report it}
\end{requirement}

\begin{requirement}{R-0085}{The surviving machine can carry the whole venue}
  A machine shall be able to run everything the venue requires of it after its
  partner is lost, and not merely its own share.
  \rationale{Where each machine is sized to carry only its own share, the surviving machine
  takes on a load it cannot carry and fails in turn. The venue then delays the
  outage rather than avoiding it, and loses the second machine as well as the
  first.}
  \uncovered{no scenario measures the survivor under the whole load}
\end{requirement}

\subsection{A partition}
\label{sec:ha:enabled:partition}

Both instances are running, and neither can see the other. Both are healthy, and
neither can establish whether the other has died or has merely become
unreachable.

This is the condition the arbiter exists for. It is the only condition in this
chapter where both available actions, promoting and declining, produce a wrong
outcome in one of the two cases.

\subsubsection{What is at risk}
\label{sec:ha:enabled:partition:risk}

An instance that cannot see its peer has two explanations available and no way to
choose between them. If it promotes itself and the peer is alive, both instances
act and their states diverge, and nothing afterwards can establish which of them
was right. If it declines and the peer has died, no instance acts. That second
failure is harder to notice, because no process crashes and nothing logs an
error, so somebody finds it by asking why trading has stopped.

\subsubsection{What the venue must do}
\label{sec:ha:enabled:partition:behaviour}

\begin{requirement}{R-0086}{An instance that cannot see its peer refers the question rather than deciding it}
  Where an instance cannot reach its peer, it shall ask the arbiter which
  instance may act, and shall not decide for itself.
  \rationale{This is the one case an instance cannot settle alone, because it is the only case
  in which a second claimant may exist that the instance cannot see. Where an
  instance can see its peer, the two settle the question between themselves under
  \reqref{R-0072}. The two rules cannot both apply at once, so no group has two
  parties entitled to decide.}
  \uncovered{no scenario partitions a pair}
\end{requirement}

\begin{requirement}{R-0087}{A partition does not produce two instances acting}
  A partition of any duration shall not result in both instances of a group
  acting.
  \rationale{\reqref{R-0058} states this property generally. This subsection restates it
  because a partition is the condition that produces the failure, and because a
  design can satisfy the general statement in every case a test produces and still
  fail here.}
  \uncovered{no scenario partitions a pair}
\end{requirement}

\subsubsection{When only one side can see the other}
\label{sec:ha:enabled:partition:asymmetric}

A partition need not be symmetric. One instance may be able to reach its peer
while the peer cannot reach it, and the two rules for settling an entitlement
then both apply at once: the instance that can see its peer settles it locally
under \reqref{R-0072}, while the instance that cannot refers it under
\reqref{R-0086}.

The two rules cannot both apply only while visibility is the same in both
directions, and here it is not. What settles this case is that an instance taking
an entitlement reports it immediately, rather than at its next scheduled report.
The arbiter therefore already holds a record of a present holder, and confirms
that holder rather than granting the entitlement again.

The window is therefore as long as a new holder takes to report itself, rather
than as long as the partition. The window is narrowed rather than removed, which
is a weaker guarantee than the venue has where one party alone issues
entitlements. A scenario should produce this case deliberately.

\subsubsection{What this costs, stated plainly}
\label{sec:ha:enabled:partition:cost}

The venue does not stop an instance on the losing side of a partition. That
instance keeps running, goes on treating itself as the leader, and is refused at
every interaction it attempts, under \reqref{R-0065}. It corrupts nothing,
because no receiver accepts what it sends. It also does not stop, and it may hold
resources until somebody intervenes.

The venue accepts that cost. Removing the instance instead would require
authority over the machine it runs on: a separate management network, hardware
power control, and credentials to use them. Those are properties of a deployment
rather than of this software, and a developer's machine has none of them, so such
a mechanism could only ever be exercised in one environment. Refusing everything
a superseded instance sends requires none of them, and is exercised wherever the
venue runs.

\subsection{Failback}
\label{sec:ha:enabled:failback}

An instance that failed is repaired and returns. This subsection settles whether
the entitlement to act returns with it.

\textbf{The entitlement stays where it is.} The instance that has just returned
is the one that recently failed, and the instance holding the entitlement is
serving the venue. Moving the position from the second to the first, on the
strength of a name fixed in configuration when the deployment was built, changes
which instance acts for a reason unconnected to either instance's fitness.

An operator may still move leadership deliberately. A machine may need to be
taken out of service, or a deployment rebalanced, and moving the position on
purpose is an ordinary operation. The venue must not move it automatically.

\begin{requirement}{R-0088}{The entitlement does not return to an instance merely because it has come back}
  A repaired instance rejoining a group shall take the follower position while
  the other instance holds the entitlement and is present, whatever its
  configured name.
  \rationale{Configured names order a cold start and decide nothing else. An automatic return
  moves the position away from an instance that is working, to the instance that
  recently failed, at a moment nobody chose.}
  \covers{ha_test:24, ha_test:32}
\end{requirement}

\begin{requirement}{R-0089}{Leadership can be moved deliberately, and moving it loses nothing}
  An operator shall be able to move the entitlement between instances of a group,
  and doing so shall lose no order, no report and no session.
  \rationale{Without this requirement, taking a machine out of service means provoking the
  failure the configuration exists to survive. An operator who must kill a healthy
  process to move a position will eventually kill the wrong one.}
  \uncovered{no scenario moves an entitlement on purpose, and nothing provides for it}
\end{requirement}

\subsection{Running degraded}
\label{sec:ha:enabled:degraded}

Every failure in this section leaves the venue in some state. For several of
them, the venue is working and can no longer survive the next failure.

A group whose second instance has not yet been restarted serves from one
instance. A venue that has lost both arbiters goes on trading, and cannot move an
entitlement if it needs to. Neither condition stops anything, neither produces an
error, and an observer outside the venue cannot distinguish either from a venue
in full health.

\textbf{The venue must go on running in these states.} Stopping when it loses
redundancy would turn the loss of a second instance into the outage that second
instance existed to prevent. The venue must also report that it is running in
such a state, because the exposure lasts until somebody restores what was lost.

\subsubsection{What the venue must do}
\label{sec:ha:enabled:degraded:behaviour}

\begin{requirement}{R-0093}{The venue trades while a group is reduced to one instance}
  The loss of one instance of a group shall not stop members trading, and shall
  not stop the surviving instance acting.
  \rationale{A pair reduced to one instance has lost its protection against the next failure,
  and has lost nothing else. Refusing to work in that state produces an outage now
  instead of risking one later.}
  \covers{ha_test:3, ha_test:6}
\end{requirement}

\begin{requirement}{R-0094}{A group running without its second instance is knowable}
  The venue shall make it possible to establish, at any moment, which groups are
  running with a second instance and which are not.
  \rationale{The period without resilience lasts from the failure until something restores the
  missing instance, and its length is the venue's exposure to a second failure.
  Automatic restart makes that period short, and also makes it invisible. Somebody
  downstream eventually notices a component that dies and stays dead, and nobody
  notices one that is restarted repeatedly. The venue must therefore report what it
  recovers automatically.}
  \uncovered{no scenario asks whether a group is running single, and nothing answers}
\end{requirement}

\begin{requirement}{R-0095}{An instance may act with no arbiter reachable, and says that it did}
  Where no arbiter can be reached and an entitlement must be settled, an instance
  may take it by the rule that orders a cold start, and shall record that it did
  so without arbitration.
  \rationale{Refusing to act because the arbiter is unreachable stops the venue for the
  failure of a component that holds no venue state and sits off the order path. It
  is also the one moment at which the venue has no protection against a second
  claimant, so the venue must record that it acted without arbitration.}
  \covers{ha_test:8, ha_test:9, ha_test:35, ha_test:39}
\end{requirement}

\subsection{How the design is checked}
\label{sec:ha:enabled:model_checking}

The scenarios in \texttt{scripts/ha\_test.py} run the venue and inject failures.
Each scenario tries one sequence of events. The failures that matter most in this
section depend on the order in which events happen: a message arriving just
after a restart, or two timeouts firing close together. A scenario that passes
says nothing about the orders it did not try.

\textbf{The design is therefore also checked by model checking.} The design is
written down in TLA+, a language for describing a system precisely enough for a
program to check it. The description states what each component holds, which
steps can change it, and which properties must always hold. The model checker,
TLC, then visits every state that a small version of the system can reach, in
every order of events, and reports any sequence of steps that breaks a property.
That sequence is called a counterexample.

The check has two limits. It covers a small version of the system: two instances,
one arbiter or two, and one or two failures at a time. It checks the design as
written in TLA+, so each property it confirms holds in the code only as far as
the TLA+ description matches the code.

\subsubsection{What is checked}
\label{sec:ha:enabled:model_checking:scope}

Two TLA+ descriptions model the design the code implements:

\begin{itemize}
  \item the two sequencer instances and the arbiter they report to; and
  \item the two arbiters and the witness, which decide which arbiter is active.
\end{itemize}

Both include crashes and restarts of every party, the failure of any link, and
messages that arrive after the state they describe has changed.

Model checking found eight defects in that design. The code now addresses six of
them in full and part of one more. The properties that now hold within the limits
checked are these:

\begin{itemize}
  \item no two leaders ever lead in the same epoch, so a receiver can always tell
        them apart;
  \item two leaders that can hear each other always resolve to one; and
  \item two active arbiters that can hear each other always resolve to one.
\end{itemize}

\textbf{Two instances can still both lead while a partition lasts.} An instance
that can reach no arbiter promotes itself, as \reqref{R-0095} permits. When the
real cause is that the instance has been cut off, the pair has two leaders until
the partition heals. The arbiter pool can likewise have two active arbiters during
a partition, because it does not yet decide by majority.

A replacement design removes promotion without an arbiter. In the replacement, an
instance leads only while a majority of three voters has granted it a lease that
has not run out. It is a proposal, described in
\texttt{docs/availability/majority\_leases.md}, and nothing in the code implements
it. A third TLA+ description specifies it. Model checking that description found
no state in which two instances act as leader at once. It also showed that each
rule of the replacement is needed: with any one rule removed, TLC finds two
instances acting as leader at once within a few steps.

The descriptions, the counterexamples and the full results are in
\texttt{docs/availability/tla/}, with the results in \texttt{findings.md}.

\subsubsection{What the venue's build must do}
\label{sec:ha:enabled:model_checking:behaviour}

\begin{requirement}{R-0145}{The build reproduces every recorded counterexample}
  The build shall rerun each counterexample recorded against a TLA+ description of
  the high availability design, and shall fail if any of them no longer breaks the
  property it is recorded against.
  \rationale{A counterexample is the evidence that a rule of the design is needed.
  A later change to the TLA+ description can stop a counterexample failing, and
  the description then no longer shows what the recorded results say it shows.
  Nothing else in the build would notice. Rerunning every counterexample takes a
  few seconds.}
  \covers{tla:lease-1-holder-counts-from-arrival, tla:lease-2-candidate-ignores-own-grant,
  tla:lease-3-restart-does-not-wait, tla:lease-4-degraded-promotion,
  tla:lease-5-epoch-regresses, tla:lease-6-no-yield-livelock, tla:lease-7-no-peer-vote}
\end{requirement}

The counterexamples reproduced today are those of the replacement design.
\texttt{scripts/tla\_trace\_pages.py} reruns them at install time, and writes a web
page that steps through each one.

